% Diagramme de classes — modèle de données principal (TikZ)
\begin{figure}[H]
\centering
\resizebox{\textwidth}{!}{%
\begin{tikzpicture}[
  class/.style={rectangle, draw=ensamBlue, fill=white, minimum width=3.0cm,
                align=left, font=\scriptsize, inner sep=4pt},
  line/.style={draw=gray!70},
  mult/.style={font=\tiny, fill=white, inner sep=1pt}
]
  \node[class] (user) at (0,0) {
    \textbf{User}\\[2pt]
    \rule{2.8cm}{0.4pt}\\[1pt]
    - id: int\\
    - name: string\\
    - email: string\\
    - role\_id: int\\
    \rule{2.8cm}{0.4pt}\\[1pt]
    + projects()\\
    + tasks()
  };
  \node[class] (role) at (4.2,0) {
    \textbf{Role}\\[2pt]
    \rule{2.8cm}{0.4pt}\\[1pt]
    - id: int\\
    - name: string\\
    - slug: string\\
    \rule{2.8cm}{0.4pt}\\[1pt]
    + permissions()
  };
  \node[class] (client) at (8.6,0) {
    \textbf{Client}\\[2pt]
    \rule{2.8cm}{0.4pt}\\[1pt]
    - id: int\\
    - name: string\\
    - email: string\\
    - status: enum\\
    \rule{2.8cm}{0.4pt}\\[1pt]
    + projects()\\
    + contacts()
  };
  \node[class] (project) at (4.2,-3.8) {
    \textbf{Project}\\[2pt]
    \rule{2.8cm}{0.4pt}\\[1pt]
    - id: int\\
    - client\_id: int\\
    - title: string\\
    - status: enum\\
    - budget: decimal\\
    \rule{2.8cm}{0.4pt}\\[1pt]
    + phases()\\
    + tasks()\\
    + quotes()
  };
  \node[class] (phase) at (0,-7.6) {
    \textbf{ProjectPhase}\\[2pt]
    \rule{2.8cm}{0.4pt}\\[1pt]
    - id: int\\
    - project\_id: int\\
    - name: string\\
    \rule{2.8cm}{0.4pt}\\[1pt]
    + tasks()
  };
  \node[class] (task) at (4.2,-7.6) {
    \textbf{Task}\\[2pt]
    \rule{2.8cm}{0.4pt}\\[1pt]
    - id: int\\
    - project\_id: int\\
    - title: string\\
    - status: enum\\
    \rule{2.8cm}{0.4pt}\\[1pt]
    + assignee()
  };
  \node[class] (quote) at (8.6,-3.8) {
    \textbf{Quote}\\[2pt]
    \rule{2.8cm}{0.4pt}\\[1pt]
    - id: int\\
    - project\_id: int\\
    - reference: string\\
    - total\_ttc: decimal\\
    \rule{2.8cm}{0.4pt}\\[1pt]
    + lines()\\
    + invoice()
  };
  \node[class] (invoice) at (12.8,-3.8) {
    \textbf{Invoice}\\[2pt]
    \rule{2.8cm}{0.4pt}\\[1pt]
    - id: int\\
    - quote\_id: int\\
    - reference: string\\
    - total\_ttc: decimal\\
  };
  \node[class] (contract) at (12.8,-7.6) {
    \textbf{Contract}\\[2pt]
    \rule{2.8cm}{0.4pt}\\[1pt]
    - id: int\\
    - project\_id: int\\
    - status: enum\\
    \rule{2.8cm}{0.4pt}\\[1pt]
    + amendments()
  };

  \draw[line] (user.east) -- node[mult, above] {N..1} (role.west);
  \draw[line] (client.south) -- ++(0,-0.5) -| node[mult, pos=0.25, left] {1} (project.north);
  \draw[line] (project.south) -- ++(0,-0.4) -| node[mult, pos=0.2, left] {1..*} (phase.north);
  \draw[line] (project.south) -- (task.north);
  \draw[line] (project.east) -- node[mult, above] {1..*} (quote.west);
  \draw[line] (quote.east) -- node[mult, above] {0..1} (invoice.west);
  \draw[line] (project.east) -- ++(1.2,0) |- node[mult, pos=0.7, right] {0..1} (contract.west);
\end{tikzpicture}%
}
\caption{Diagramme de classes simplifié du modèle de données métier}
\label{fig:class_diagram}
\end{figure}
